\clearpage\label{answers}
{\sans\bfseries\LARGE\color{ink} Answers}\par
\footnotesize\linespread{0.94}\selectfont
Try every puzzle first. Where an answer shows a calculation, check the steps, not just the result.

\subhead{Chapter 1}
\textbf{Multix on vanishing.} No. The covered cup kept its water, so the water in the open cup most likely moved into the air. \textbf{Keep box.} Kept: how much water was left. Set aside: the path the water took.

\textbf{1.1} $40-35=5$~mm. 0~mm.\\
\textbf{1.2} Drops of 3, 2, 3 and 2~mm, a total of 10~mm in four days, so $10\div4=2.5$~mm a day.\\
\textbf{1.3} $10\times2=20$, so $50-20=30$~mm. It is a guess because the rate may change with weather.\\
\textbf{1.4} No. The reading went to 59 and back to 60, so it is a wobble.\\
\textbf{1.5} Two things changed, the lid and the place (sun against dark).\\
\textbf{1.6} It makes ``stopped existing'' and ``soaked away'' less likely, since the covered cup in the same glass kept all its water. The water moved into the air.\\
\textbf{1.7} Air gained 11 and the cup now holds 49.\\
\textbf{1.8} The sun, which warms the water, and perhaps a breeze. The shade is cooler.\\
\textbf{1.9} Answers vary, such as wet washing, a wet towel or a puddle.\\
\textbf{1.10} The covered cup would lose water too. It did not.

\subhead{Chapter 2}
\textbf{Multix on cold days.} No. Evaporation is faster in the warm, but it also happens in the cold. Washing dries slowly on a cold windy day. \textbf{Keep box.} Kept: which dried sooner. Set aside: the path of the vapour.

\textbf{2.1} Flat: $36+40+38=114$, average 38. Folded: $120+126+132=378$, average 126.\\
\textbf{2.2} Line in a breeze 1, heap 2, folded in a closed drawer 3.\\
\textbf{2.3} The size of the square and the place (sun against shade).\\
\textbf{2.4} $135\div42\approx3.2$, so about 3 times.\\
\textbf{2.5} (a) D (b) V (c) D.\\
\textbf{2.6} Washing still dries on a cold day, so evaporation is not only for hot weather.\\
\textbf{2.7} The air gained 4 spoonfuls.\\
\textbf{2.8} In a fold less of the wet paper touches open air, so the vapour escapes more slowly.\\
\textbf{2.9} The wide tray. The water is spread over more area.\\
\textbf{2.10} Answers vary. Check that only one thing is changed.

\subhead{Chapter 3}
\textbf{Multix on the leak.} It missed the colour test. Coloured water inside gave clear drops outside, which points to the air. \textbf{Keep box.} Kept: water moved from the air to the glass. Set aside: the number and size of the drops.

\textbf{3.1} $90+80+100=270$, and $270\div3=90$ seconds.\\
\textbf{3.2} Came from the air.\\
\textbf{3.3} It was not cold enough to cool the air next to it, so no vapour turned into drops.\\
\textbf{3.4} The temperature of the glass and the place (kitchen against outside in the sun).\\
\textbf{3.5} Examples: a cold can, a window on a cold morning, a cold spoon.\\
\textbf{3.6} Air after 48, glass after 2.\\
\textbf{3.7} (a) S (b) G (c) B.\\
\textbf{3.8} The window is cold. The vapour comes from the warm air in the room.\\
\textbf{3.9} The film is cooler than the vapour, so the vapour turns into drops on it.

\subhead{Chapter 4}
\textbf{Multix on the missing spoonful.} Not yet. A gap of one spoonful is about the size of the counting error, and some water is on the lid and sides. \textbf{Keep box.} Kept: the total. Set aside: the exact path.

\textbf{4.1} $40-31=9$.\\
\textbf{4.2} $38+2=40$. Yes.\\
\textbf{4.3} $36+3=39$, so the gap is 1. It could be a counting error. Recount or repeat.\\
\textbf{4.4} Drops run down 4, water evaporates 1, drops form on the lid 3, vapour reaches the lid 2.\\
\textbf{4.5} The water goes round inside the jar, so the soil is watered again.\\
\textbf{4.6} Jar A was closed and Jar B was open.\\
\textbf{4.7} ``8 spoonfuls moved into the air as vapour.''\\
\textbf{4.8} 18 stays, and the air gained 6.\\
\textbf{4.9} A sealed jar keeps the vapour in, so the total can be counted.\\
\textbf{4.10} Answers vary.

\subhead{Chapter 5}
\textbf{Multix on new water.} No. The clear water in the cup came from the bowl by evaporation and condensation. \textbf{Keep box.} Kept: the order of the stages. Set aside: sizes, places and times.

\textbf{5.1} $20-15-4=1$. It is probably drops on the film or the sides.\\
\textbf{5.2} Evaporates 1, vapour rises 2, condenses into cloud drops 3, rain falls 4, rivers flow to the sea 5.\\
\textbf{5.3} (a) E (b) C (c) R (d) E (e) C (f) R.\\
\textbf{5.4} The water travelled as vapour. The colour stayed behind.\\
\textbf{5.5} The first route.\\
\textbf{5.6} Snow or hail.\\
\textbf{5.7} Examples: plants, ice, underground water, the amounts, the times.\\
\textbf{5.8} The real steps are not the same size or speed.\\
\textbf{5.9} Answers vary. Water can stay in the deep ocean or in ice for a very long time.\\
\textbf{5.10} Puddle to air to cloud to rain and back to the puddle.

\subhead{Chapter 6}
\textbf{Multix on the diagram.} No. A five-box diagram keeps the order but not amounts or times. \textbf{Keep box.} Kept: the total rain. Set aside: which day it fell.

\textbf{6.1} $5+0+3+6+1=15$~mm, average 3~mm.\\
\textbf{6.2} (a) Y (b) N (c) N.\\
\textbf{6.3} Answers vary, for example numbers on the arrows to say how much.\\
\textbf{6.4} (a) E (b) C (c) R.\\
\textbf{6.5} ``The water in the open cup moved into the air.''\\
\textbf{6.6} $30-22=8$.\\
\textbf{6.7} Trees catch the rain, so two things changed.\\
\textbf{6.8} (a) M (b) R (c) M.\\
\textbf{6.9} Example: how much water is left in a cup after five days?\\
\textbf{6.10} $8+0+0+4+3=15$~mm.\\
\textbf{6.11} Evaporation, then condensation.\\
\textbf{6.12} Condensation.\\
\textbf{6.13} One jar does not show water is never lost anywhere. Check other jars and conditions.\\
\textbf{6.14} The lake holds 94 and the air 26.\\
\textbf{6.15, 6.16} Answers vary. Check that water moves and nothing vanishes.\\
\textbf{Slop check.} (a) You cannot tell: no title, no labels. (b) Many colours, stars, joke text. (c) A title, the days and the rain amounts with units. (d) The stars and jokes.

\clearpage\label{notes}
{\sans\bfseries\LARGE\color{ink} Words}\par
\small
\begin{multicols}{2}
\word{condensation}{vapour turning back into drops of liquid.}
\word{evaporation}{liquid turning into vapour.}
\word{ledger}{a list of where the amounts are, so a total can be added up.}
\word{model}{a simple stand-in that keeps some things and drops others.}
\word{terrarium}{a sealed jar with soil and a plant.}
\word{vapour}{water as an invisible gas in the air.}
\word{water cycle}{the loop of water from sea to air to rain and back.}
\word{spread}{how far apart repeated readings are.}
\end{multicols}

\vspace{0.4em}
{\sans\bfseries\large\color{sky!70!black} Notes for grown-ups}\par
This book teaches that a diagram of the water cycle is a model: it keeps the order of the stages and drops the times, places and amounts. Each chapter has one activity, one worked example with made-up practice numbers, one trust box, one ``what did we keep?'' prompt and a page of puzzles.

\textbf{Materials and safety.} Clear cups and jars, a ruler, cling film, ice, food colouring, a zip bag. Use cold tap water only. Do not drink water from an experiment. Clean up spills.

\textbf{Vapour and steam.} Water vapour is invisible. Mist, steam clouds and clouds in the sky are tiny drops of liquid water.

\textbf{The practice numbers} were made up so the arithmetic is easy to check. Real results will vary with the room, the day and the equipment.

\subhead{Sources}
\begin{itemize}[leftmargin=1.3em,itemsep=0.2em]
\item U.S. Geological Survey, Water Science School. \emph{The Water Cycle}. For the stages and how long water stays in the air.
\item C. D. Ahrens. \emph{Meteorology Today}. Cengage. For condensation, dew and clouds.
\item J. R. Taylor. \emph{An Introduction to Error Analysis}, 2nd edition. University Science Books, 1997. For repeats and spread.
\end{itemize}

\clearpage
\thispagestyle{blank}
\begin{tikzpicture}[remember picture,overlay]
  \fill[sun!35] ($(current page.south west)+(0.4in,0.4in)$) rectangle ($(current page.north east)-(0.4in,0.4in)$);
  \node[anchor=north,align=center,font=\sffamily\bfseries,text=ink,text width=5.4in] at ($(current page.north)+(0,-1.4in)$) {\Large Next in Science\\[6pt]\fontsize{28}{30}\selectfont Mixed Up and Sorted Out};
  \node[align=center,font=\sffamily,text width=4.6in,text=ink] at (current page.center) {A spoonful of salt goes into water and seems to vanish. Book 6 shows how to separate mixtures, and how to check that nothing was lost.};
  \node at ($(current page.south)+(0,2.6in)$) {\scalebox{3}{\multixbot[ok]}};
  \node[font=\sffamily\bfseries,text=ink] at ($(current page.south)+(0,0.9in)$) {Ages 9 to 11 \quad$\star$\quad Science Book 5};
\end{tikzpicture}
